\clearpage
\VerdierPhysicalPageAnchor{0149}%
\begingroup
\small
\setlength{\parskip}{0pt}%
\setlength{\emergencystretch}{1.5em}%
\setlength{\abovedisplayskip}{3pt plus 1pt minus 1pt}%
\setlength{\belowdisplayskip}{3pt plus 1pt minus 1pt}%

\VerdierSemanticAnchor{ch2-2-3-5-injectivity-p149-continuation}
\noindent 使得 $uf=0$。再用 a), i)，可知存在交换图：
\VerdierSemanticAnchor{ch2-2-3-5-injectivity-refinement-diagram}
\begin{center}
 \begin{tikzcd}[column sep=3.7em,row sep=2.3em]
 && {} \\
 & {} \arrow[ur,"v"] && {t'\in S_{\mathcal A}(\mathcal A\cap\mathcal B)} \\
 X \arrow[ur,"f"] && Y \arrow[ul,"s"] \arrow[uu,"{t'}"'] &
 \end{tikzcd}
\end{center}
\VerdierSemanticAnchor{ch2-2-3-5-injectivity-conclusion}
\noindent 使得 $vf=0$。于是 $m=Q_1(t')^{-1}Q_1(vf)=0$。
\VerdierSemanticAnchor{ch2-2-3-5-surjectivity-map}
现在证明同态
\[
 \operatorname{Hom}_{\mathcal A/(\mathcal A\cap\mathcal B)}(X,Y)
 \longrightarrow\operatorname{Hom}_{\mathcal D/\mathcal B}(X,Y)
\]
是满射。
\VerdierSemanticAnchor{ch2-2-3-5-surjectivity-roof}
任意元素 $p\in\operatorname{Hom}_{\mathcal D/\mathcal B}(X,Y)$ 可写成
\[
 p=Q_2(t)^{-1}Q_2(g),\qquad t\in S_{\mathcal D}(\mathcal B).
\]
\VerdierSemanticAnchor{ch2-2-3-5-surjectivity-refinement}
利用 a), i)，这种元素总可写成
\[
 p=Q_2(s)^{-1}Q_2(g'),\qquad
 s\in S_{\mathcal A}(\mathcal A\cap\mathcal B),
\]
\VerdierSemanticAnchor{ch2-2-3-5-surjectivity-conclusion}
\noindent 因而来自
$\operatorname{Hom}_{\mathcal A/(\mathcal A\cap\mathcal B)}(X,Y)$ 的元素。
\VerdierSemanticAnchor{ch2-2-3-5-c-proof-i-implies-ii}
还需证明 c) 中的等价性。先证明 i) $\Rightarrow$ ii)。
\VerdierSemanticAnchor{ch2-2-3-5-iterated-quotient}
由 (2.3.1)，范畴 $\mathcal D/\mathcal B$ 典范同构于
$(\mathcal D/(\mathcal A\cap\mathcal B))
/(\mathcal B/(\mathcal A\cap\mathcal B))$，且典范函子
\VerdierSemanticAnchor{ch2-2-3-5-q-functor}
\[
 Q\colon\mathcal D/(\mathcal A\cap\mathcal B)\longrightarrow\mathcal D/\mathcal B
\]
是商函子。
\VerdierSemanticAnchor{ch2-2-3-5-q-free-reduction}
因此由命题 (2.3.3, a))，为证明 i) $\Rightarrow$ ii)，只需证明：
对 $\mathcal A$ 的每个对象 $X$ 与
$\mathcal B$ 的每个对象 $B$，
$\operatorname{Hom}_{\mathcal D/(\mathcal A\cap\mathcal B)}(B,X)=0$。
\VerdierSemanticAnchor{ch2-2-3-5-m-representation}
$\operatorname{Hom}_{\mathcal D/(\mathcal A\cap\mathcal B)}(B,X)$ 的任意元素
$m$ 可写成 $Q_3(f)Q_3(s)^{-1}$，其中
\VerdierSemanticAnchor{ch2-2-3-5-q3-functor}
\[
 Q_3\colon\mathcal D\longrightarrow\mathcal D/(\mathcal A\cap\mathcal B)
\]
是典范函子，且 $s$ 属于 $S_{\mathcal D}(\mathcal A\cap\mathcal B)$；后者是
$\mathcal D$ 中与 $\mathcal A\cap\mathcal B$ 关联的乘法系 (2.2.2)。
于是 $\mathcal D$ 中有图式：
\VerdierSemanticAnchor{ch2-2-3-5-w-diagram}
\begin{center}
 \begin{tikzcd}[column sep=4em,row sep=2.4em]
 & W \arrow[dl,"s"'] \arrow[dr,"f"] & \\
 B && X
 \end{tikzcd}\quad ,
\end{center}
\endgroup
